% Introduction and Related Work drafts -- to be integrated into main.tex

\section{Introduction}

The success of self-supervised foundation models in natural language processing has catalysed a new generation of DNA language models that aim to learn transferable representations of genomic sequence without task-specific supervision. Building on the transformer architecture \citep{vaswani2017attention,devlin2019bert} and recent efficient alternatives such as Hyena operators \citep{poli2023hyena} and selective state-space models \citep{gu2023mamba}, several DNA foundation models have appeared in rapid succession. DNABERT-2 extends the original DNABERT line with byte pair encoding and ALiBi positional embeddings trained on 135 species \citep{zhou2024dnabert2}; the Nucleotide Transformer family (NT-v2) was trained on more than 850 genomes and demonstrated cross-species transfer for regulatory tasks \citep{dallatorre2025nt}; HyenaDNA replaces attention with gated long convolutions and scales to million-base contexts on the human reference genome \citep{nguyen2023hyenadna}; Caduceus-Ph builds on MambaDNA with reverse-complement equivariance \citep{schiff2024caduceus}; and GROVER tokenises the human genome with a learned byte pair encoding \citep{sanabria2024grover}. Alongside these general models, specialised supervised predictors such as DeepSEA \citep{zhou2015deepsea}, Basset \citep{kelley2016basset}, DanQ \citep{quang2016danq}, SpliceAI \citep{jaganathan2019spliceai}, Enformer \citep{avsec2021enformer} and Sei \citep{chen2022sei} have continued to set strong baselines on specific regulatory and expression endpoints.

Despite this progress, quantitative comparisons between DNA foundation models remain difficult to interpret. Published benchmarks are almost always conducted after task-specific fine-tuning, which confounds the representational capacity of a frozen model with idiosyncratic choices of hyperparameters, parameter-efficient fine-tuning methods \citep{yuan2026borzoi_peft}, and the particular downstream layer on top of the backbone. A recent zero-shot approach that appends a small convolutional head to frozen embeddings partially mitigates this confound, but it was restricted to a handful of human-centric tasks and left pooling strategy -- how sequence-level predictions are aggregated from token embeddings -- largely unexplored. Moreover, current benchmarks rarely span the full diversity of questions that genome-scale models are expected to answer, from functional element classification and epigenetic modification prediction to gene expression quantification \citep{avsec2021enformer,karollus2023enformer_limits}, variant effect quantification for pathogenic alleles \citep{landrum2018clinvar,jaganathan2019spliceai} and fine-mapped QTLs \citep{gtex2020}, and the recognition of higher-order chromatin architecture such as topologically associating domains \citep{dixon2012topological,lieberman2009hic}.

Here we present a zero-shot embedding benchmark designed to close these gaps. We evaluate five representative DNA foundation models -- DNABERT-2, NT-v2, HyenaDNA, Caduceus-Ph and GROVER -- on 57 sequence classification datasets (52 binary, 5 multi-class) spanning human and multi-species genome region classification, human and multi-species epigenetic modification, and multi-class regulatory classification. We train standard, lightly tuned random forest classifiers (with naive Bayes and elastic-net logistic regression as sensitivity checks and a baseline CNN) on top of frozen embeddings produced by three pooling strategies: sentence-level summary tokens, mean token pooling and maximum pooling. We then extend the benchmark beyond classification to gene expression prediction from paired sequence/expression data (GTEx v8 whole blood, 610 subjects, 21{,}004 genes at 6 kb and 768 genes at 196 kb); to variant effect quantification for 22{,}222 pathogenic and 17{,}374 common SNPs via paired reference/alternative sequences; to QTL classification across 1{,}896 eQTLs, 540 sQTLs, 116 ipaQTLs and 142 paQTLs from GTEx v8 fine-mapped fine-mapping; to a controlled pre-training experiment in which HyenaDNA is re-pretrained on DNABERT-2's multi-species corpus (135 species, 6 taxonomic categories); and to a TAD-attention analysis on 1{,}500 TAD-centered versus 1{,}500 background 6{,}000 bp sequences.

Our contributions are as follows:
\begin{itemize}
  \item We demonstrate that mean token pooling consistently and significantly outperforms summary-token and maximum pooling across all five models and 52 binary classification datasets, with average AUC gains of 1.4\%--8.7\% (DeLong's test \citep{delong1988comparing}, $p<0.01$).
  \item We show that general DNA foundation models can surpass specialised supervised predictors on pathogenic variant identification (NT-v2 AUC 0.73, Cohen's $d=0.88$; Enformer AUC 0.69, $d=0.73$), while specialised models dominate QTL prediction (e.g.\ AlphaGenome sQTL AUC 0.80, ipaQTL AUC 0.86) \citep{avsec2021enformer,chen2022sei,avsec2025alphagenome}.
  \item We provide controlled evidence that multi-species pre-training improves 14 of 49 datasets for HyenaDNA, including human 5mC detection (0.707$\to$0.749) and human vs.\ worm classification (0.968$\to$0.984).
  \item We report that NT-v2 self-attention averaged over all layers and heads does not recognise TAD boundaries without fine-tuning, marking a clear boundary of current zero-shot capability.
  \item We release a uniform, zero-shot evaluation protocol that enables rapid, pooling-controlled comparison of future DNA foundation models.
\end{itemize}

\section{Related Work}

\paragraph{Supervised sequence-to-function models.}
Early convolutional networks such as DeepBind \citep{alipanahi2015deepbind}, DeepSEA \citep{zhou2015deepsea}, Basset \citep{kelley2016basset} and DanQ \citep{quang2016danq} established that regulatory activity and variant effects can be predicted directly from short DNA sequences. Enformer \citep{avsec2021enformer} scaled the receptive field to 100 kb using transformer attention, and Sei \citep{chen2022sei} learned a global map of sequence classes by predicting 21{,}907 chromatin profiles across more than 1{,}300 cell lines and tissues. SpliceAI \citep{jaganathan2019spliceai} achieved near-experimental accuracy on splice site prediction and quantified cryptic splicing variants. These models remain strong baselines for specific endpoints, but they require ground-truth regulatory labels for training and do not offer the task-agnostic representations that motivate DNA foundation models. Recent work has also shown that even powerful supervised predictors such as Enformer largely capture promoter-proximal signal and struggle to integrate causal effects of distal enhancers \citep{karollus2023enformer_limits}.

\paragraph{DNA foundation models.}
Self-supervised DNA language models pre-train on unlabelled genomic sequence to produce reusable embeddings. DNABERT-2 trains a BERT backbone with byte pair encoding on 135 species \citep{zhou2024dnabert2}; the Nucleotide Transformer family extends this to 850 genomes with rotary positional embeddings and up to 2.5 billion parameters \citep{dallatorre2025nt}. HyenaDNA replaces attention with implicit long convolutions (Hyena operators \citep{poli2023hyena}) and scales to million-base contexts on the human reference genome \citep{nguyen2023hyenadna}. Caduceus-Ph adapts the Mamba selective state-space model \citep{gu2023mamba} with reverse-complement equivariance for DNA \citep{schiff2024caduceus}, and GROVER tokenises the human genome via a learned byte pair vocabulary \citep{sanabria2024grover}. Downstream studies have begun to exercise these models in realistic genomic settings -- for example, on polyadenylation site prediction \citep{saha2025polya_glm}, plant DNA methylation \citep{feng2025hyenadna_methylation}, gene-fusion breakpoint classification \citep{krupicka2026fusion}, and Arabidopsis variant effect scoring \citep{benegas2023gpn} -- yet each focuses on a narrow task and typically uses a single pooling strategy.

\paragraph{Variant effect and QTL prediction.}
Variant effect prediction is a central application for genomic models. Pathogenic variants curated in ClinVar \citep{landrum2018clinvar} provide a reference benchmark distinct from expression quantitative trait loci derived from fine-mapped GTEx v8 whole-blood data \citep{gtex2020}. Specialised models such as Enformer, Sei and AlphaGenome \citep{avsec2021enformer,chen2022sei,avsec2025alphagenome} have produced state-of-the-art scores on QTL prediction tasks, and parameter-efficient fine-tuning of Borzoi further extends supervised pipelines to novel cellular contexts \citep{yuan2026borzoi_peft}. Whether general DNA foundation models offer competitive zero-shot representations for variant effect tasks, and how this depends on variant class and effect size, has not been systematically quantified.

\paragraph{Chromatin structure and higher-order organisation.}
Topologically associating domains (TADs) are megabase-scale self-interacting chromatin units conserved across cell types and species \citep{dixon2012topological,lieberman2009hic}. Whether attention maps from zero-shot DNA language models recognise TAD boundaries is an open question that directly probes how much higher-order information is available in raw sequence patterns.

\paragraph{Positioning of this work.}
Our study differs from prior DNA benchmarks in three respects. First, we explicitly treat pooling as a first-class hyperparameter and report its effect on every classification task. Second, we combine classification, expression regression, variant effect and chromatin structure probes in a single framework so that task-specific strengths and weaknesses are directly comparable. Third, we include a controlled pre-training data swap for HyenaDNA to isolate the contribution of multi-species context independently of architectural differences.
